\chapter{The Part I systems view}
\label{ch:part-i-overview}

\begin{learningobjectives}
After this chapter, you should be able to:
\begin{itemize}
  \item distinguish a learned model from the data-driven system that contains it;
  \item trace a result across data, computational, software, statistical,
  scientific, and operational boundaries;
  \item state what each Part I chapter contributes to that system; and
  \item match a correctness claim with evidence of the appropriate scope.
\end{itemize}
\end{learningobjectives}

\begin{chapterconnection}
This chapter fixes the level of abstraction for Part I. Python objects,
statistical models, databases, APIs, and deployed applications will be studied
as components of one evidence-bearing system. Later chapters supply the
technical details; this chapter states the contracts those details must serve.
\end{chapterconnection}

\section{The model is one component}

A machine-learning project is often summarized by its model: a random forest,
neural network, or large language model. That description omits most of the
system. A model does not determine whether measurements have the correct units,
whether labels represent the intended quantity, whether features will exist at
prediction time, or whether a result is safe to use. It also does not version
its inputs, expose an interface, monitor its behavior, or recover from failure.

\begin{definitionbox}{Machine-learning model}
A \emph{machine-learning model} is a parameterized rule whose behavior is
partly determined from data. Given an input representation, it produces an
output such as a score, category, estimate, generated object, or proposed
action.
\end{definitionbox}

\begin{definitionbox}{Data-driven software system}
A \emph{data-driven software system} is a collection of data, programs,
interfaces, infrastructure, and operating procedures whose behavior depends
materially on observed data. It transforms evidence into a prediction,
decision, explanation, or scientific claim under explicit technical and human
responsibilities.
\end{definitionbox}

If a model is written as
\[
  f_{\theta}: \mathcal{X} \longrightarrow \mathcal{Y},
\]
then model fitting determines parameters \(\theta\) from training data. An
operated prediction requires a larger composition:
\[
  r
  \xrightarrow{\text{parse}}
  v
  \xrightarrow{\text{validate}}
  x
  \xrightarrow{f_{\theta}}
  \widehat{y}
  \xrightarrow{\text{present}}
  a.
\]
Here \(r\) is a raw record, \(v\) a typed value, \(x\) the model input,
\(\widehat{y}\) the model output, and \(a\) the result presented to a person or
another program. A defect in any arrow can invalidate the final result even
when \(f_{\theta}\) behaves exactly as tested.

\begin{definitionbox}{Engineering boundary}
An \emph{engineering boundary} is a point at which data, control, ownership, or
trust passes between contexts. A function call is a small boundary; a file,
database, remote API, deployment, and human approval cross larger ones. Each
boundary requires a contract: what enters, what leaves, what may fail, and who
is responsible.
\end{definitionbox}

\subsection{Three views of the same work}

The same computation must be legible from three views:

\begin{center}
\small
\begin{tabularx}{\textwidth}{@{}p{0.20\textwidth}X X@{}}
\toprule
\textbf{View} & \textbf{Central questions} & \textbf{Typical failure} \\
\midrule
Scientific & What question is being asked? What observations and assumptions
support the claim? & A proxy is treated as the phenomenon of interest, or
prediction is mistaken for causation. \\
Computational & How is meaning represented and transformed? What numerical,
time, and memory limits apply? & Missingness becomes zero, an array has the
wrong shape, or an algorithm is numerically unstable. \\
Operational & How does behavior reach a user? Who may invoke it, how is it
observed, and how is it recovered? & A stale model is served, an interface
changes meaning, or a failure has no owner. \\
\bottomrule
\end{tabularx}
\end{center}

For a classifier predicting material stability, the scientific view asks
whether the experiment measures stability and whether the evaluation supports
the stated population. The computational view asks how compositions,
temperatures, and missing values become features. The operational view asks
whether the service uses the same feature definition and identifies inputs
outside the model's domain. These are not separate projects; they are separate
claims about one result.

\begin{misconceptionbox}
\textbf{Misconception:} a notebook that runs from top to bottom and reports high
test accuracy constitutes a complete machine-learning system.

\textbf{Correction:} successful execution describes one run in one
environment. Test accuracy describes one metric under a sampling design.
Neither alone establishes valid data, reproducibility, deployability, safety,
or scientific relevance.
\end{misconceptionbox}

\subsection{The evidence loop}

An operated system is a loop rather than a pipeline that ends at prediction.

\begin{center}
\begin{tikzpicture}[node distance=5mm and 5mm]
  \node[flowbox] (evidence) {Evidence\\files, instruments};
  \node[flowbox,right=of evidence] (meaning) {Typed meaning\\schema, units};
  \node[flowbox,right=of meaning] (analysis) {Analysis\\features, models};
  \node[flowbox,below=of analysis] (product) {Data product\\API, client};
  \node[flowbox,left=of product] (operate) {Operation\\release, monitoring};
  \node[flowbox,left=of operate] (learn) {Revision\\tests, redesign};
  \draw[flowarrow] (evidence) -- (meaning);
  \draw[flowarrow] (meaning) -- (analysis);
  \draw[flowarrow] (analysis) -- (product);
  \draw[flowarrow] (product) -- (operate);
  \draw[flowarrow] (operate) -- (learn);
  \draw[flowarrow] (learn) -- (evidence);
\end{tikzpicture}
\end{center}

At each arrow, meaning changes representation or context. A file reader needs a
format, encoding, schema, and missing-value policy. A prediction endpoint needs
an input schema, authentication rule, model identity, output definition, and
failure response. After release, logs, outcomes, and user reports become new
evidence. They may justify a test, a retrained model, a changed interface, or
retirement of the system.

\begin{workedexample}{A plausible answer from a broken system}
An instrument records \code{-999} when a temperature sensor fails. A program
parses the column as numbers, converts every value from Fahrenheit to Celsius,
and sends the result to a well-evaluated model. The model returns \(0.82\), and
the interface displays \emph{82\% likely stable}.

The conversion and model may both satisfy their local tests. The composition is
wrong because a sentinel was treated as a physical measurement. A defensible
boundary maps the documented sentinel to missingness, checks the physical
range, preserves the raw record, and records whether the observation was
rejected or imputed. Representation is part of the scientific argument, not
merely data cleaning.
\end{workedexample}

\subsection{Why learned components require additional controls}

Ordinary software already requires clear interfaces, version control, tests,
security, and observability. Learned behavior adds four complications:

\begin{enumerate}
  \item \textbf{Data changes behavior.} Training population, labels,
  preprocessing, and sampling can alter a model without changing its algorithm.
  \item \textbf{Evaluation is conditional.} A metric has meaning only with its
  dataset, split design, baseline, uncertainty, and intended use.
  \item \textbf{Behavior can drift.} Instruments, schemas, populations, and
  user behavior change after release.
  \item \textbf{Outputs are not authority.} A score or agent proposal does not
  acquire permission to create a consequential side effect merely because it
  was produced automatically.
\end{enumerate}

Scientific discovery raises the standard further. Prediction may support a
hypothesis, parameter estimate, mechanism, or experimental decision, but it is
not identical to any of them. A scientific system must distinguish observation
from inference, association from intervention, exploration from confirmation,
and model uncertainty from measurement uncertainty.

\begin{definitionbox}{Data product}
A \emph{data product} is a maintained interface through which people or
software obtain value from data. It may be a dataset, report, dashboard, API,
decision-support tool, scientific workflow, or automated service. A model may
be one component; contracts, users, delivery, operation, and feedback complete
the product.
\end{definitionbox}

The appropriate architecture is the smallest one that protects the required
contracts. A CSV file can be correct for a small immutable dataset. A database
is justified by durable identity, queries, transactions, concurrency, or scale.
An API is justified when a responsibility must cross a process, machine,
language, team, or trust boundary.

\conceptcheck{For a machine-learning application you know, identify the model,
the data product, one scientific assumption, one software boundary, and the
role responsible for recovery.}

\section{The chapter map}

Part I expands the engineering boundary while retaining earlier contracts. A
deployed API still depends on correctly interpreted strings; a tool-using agent
still depends on ordinary functions, exceptions, credentials, and tests.

\begin{center}
\begin{tikzpicture}[node distance=7mm and 5mm]
  \node[flowbox,text width=27mm] (environment) {1. Environment\\interpreter, shell};
  \node[flowbox,text width=27mm,right=of environment] (objects) {2. Objects\\values, collections};
  \node[flowbox,text width=27mm,right=of objects] (contracts) {3. Contracts\\functions, classes, files};
  \node[flowbox,text width=27mm,right=of contracts] (project) {4. Project\\package, tests, CI};
  \node[flowbox,text width=27mm,below=of project] (evidence) {5. Evidence\\arrays, tables, plots};
  \node[flowbox,text width=27mm,left=of evidence] (persistence) {6. Persistence\\SQL, remote data};
  \node[flowbox,text width=27mm,left=of persistence] (agency) {7. Agency\\models, tools, policy};
  \node[flowbox,text width=27mm,left=of agency] (product) {8. Product\\API, client, operations};
  \draw[flowarrow] (environment) -- (objects);
  \draw[flowarrow] (objects) -- (contracts);
  \draw[flowarrow] (contracts) -- (project);
  \draw[flowarrow] (project) -- (evidence);
  \draw[flowarrow] (evidence) -- (persistence);
  \draw[flowarrow] (persistence) -- (agency);
  \draw[flowarrow] (agency) -- (product);
\end{tikzpicture}
\end{center}

\subsection{Chapter-by-chapter capabilities}

\begin{longtable}{@{}
  >{\raggedright\arraybackslash}p{0.05\textwidth}
  >{\raggedright\arraybackslash}p{0.20\textwidth}
  >{\raggedright\arraybackslash}p{0.32\textwidth}
  >{\raggedright\arraybackslash}p{0.34\textwidth}@{}}
\toprule
\textbf{Ch.} & \textbf{Boundary} & \textbf{Result} &
\textbf{Evidence} \\
\midrule
\endhead
1 & Computing environment & Project-local Python environment and matching VS
Code kernel & Interpreter path, version, and reproducible setup check \\
2 & Native objects & Transformations using strings, sequences, mappings, sets,
and iteration & Predicted values and types, explicit mutation, edge cases \\
3 & Functions, classes, files & Documented domain model and native-file
knowledge graph & Type hints, NumPy-style documentation, exceptions,
round-trip tests \\
4 & Package and workflow & Installable, versioned package with CI & Unit,
integration, doctest, and notebook checks in a clean environment \\
5 & Numerical evidence & Array- and table-based analysis with visual evidence &
Shape and dtype checks, labels, justified encodings, reproducible figures \\
6 & Durable data & Parameterized query and memory-bounded workflow & Schema,
transaction, query-safety, and chunking evidence \\
7 & Model and agent & Provider-independent tool workflow with bounded actions &
Secret isolation, structured responses, evaluations, budgets, approvals \\
8 & Operated product & Value traced from storage through API to client &
Contract tests, release checks, monitoring, and recovery plan \\
\bottomrule
\end{longtable}

The sequence has four dependencies. Chapters 1--2 establish computational
identity and Python's object model. Chapters 3--4 turn reasoning into documented,
tested, reusable software. Chapters 5--6 introduce numerical structures and
durable data access. Chapters 7--8 place probabilistic and automated components
inside operated systems.

\subsection{Milestones}

\begin{description}
  \item[After Chapter 2: reproducible computation.] A reader can select the
  correct interpreter, run the work, inspect its objects, and explain its native
  transformations.
  \item[After Chapter 4: maintainable component.] Reusable logic lives in an
  installable package with documented interfaces, tests, and a version.
  \item[After Chapter 6: evidence workflow.] The project retrieves, validates,
  transforms, and visualizes data without assuming all data fits in memory.
  \item[After Chapter 8: operated product.] The project exposes a bounded,
  observable interface with explicit release and recovery procedures.
\end{description}

These are cumulative capabilities. When a later boundary exposes a weak earlier
contract, the earlier contract must be repaired.

\conceptcheck{Choose one chapter-map row. Explain why its result is insufficient
without the listed evidence, and identify one later chapter that depends on it.}

\section{Correctness has layers}

The assertion \emph{the result is correct} is incomplete unless it names a
requirement, context, and evidence. A program may compute a formula correctly
while the formula answers the wrong scientific question. A model may be
statistically sound while a service feeds it a different feature
representation.

\begin{definitionbox}{Correctness claim}
A \emph{correctness claim} states that an artifact or behavior satisfies a
specified requirement under stated conditions. \emph{The model works} is weak.
\emph{Model 1.4 meets the prespecified recall threshold on the held-out 2026
instrument dataset using feature pipeline 3} is testable.
\end{definitionbox}

\subsection{Six lenses for evaluating a result}

\begin{longtable}{@{}
  >{\raggedright\arraybackslash}p{0.15\textwidth}
  >{\raggedright\arraybackslash}p{0.28\textwidth}
  >{\raggedright\arraybackslash}p{0.27\textwidth}
  >{\raggedright\arraybackslash}p{0.22\textwidth}@{}}
\toprule
\textbf{Lens} & \textbf{Question} & \textbf{Representative evidence} &
\textbf{Limit} \\
\midrule
\endhead
Data & Do values represent the correct entities, units, identity, and
missingness? & Provenance, schema, range, calibration, and duplicate checks &
Does not establish that the measurements answer the question \\
Computational & Does the algorithm implement the specified mathematics within
acceptable error? & Reference cases, invariants, convergence and tolerance
analysis & Does not establish that the algorithm is appropriate \\
Software & Do components satisfy their contracts when composed? & Unit,
integration, contract, security, and dependency checks & Tested cases may not
represent future operation \\
Statistical & Does the design support the estimator or generalization claim? &
Sampling rationale, baselines, held-out evaluation, calibration, intervals &
Association does not imply intervention or importance \\
Scientific & Do evidence and assumptions support the domain claim? & Controls,
theory, experimental design, replication, alternative explanations & Does not
establish that deployment preserves the analysis \\
Operational & Does the released system deliver intended behavior in use? &
Version records, logs, metrics, traces, audits, outcomes & Reliability does not
establish scientific validity \\
\bottomrule
\end{longtable}

The limit column prevents evidence from being asked to prove too much. A unit
test can verify a Fahrenheit-to-Celsius function; it cannot establish that the
input was Fahrenheit. A held-out set estimates behavior under a sampling
assumption; it cannot guarantee an unchanged deployment population.

\subsection{Verification, validation, and reproducibility}

\begin{definitionbox}{Verification}
\emph{Verification} asks whether an implementation satisfies its specification:
did we build the specified thing correctly? Tests, type checks, static analysis,
schema checks, and comparison with known results provide verification evidence.
\end{definitionbox}

\begin{definitionbox}{Validation}
\emph{Validation} asks whether the specification and system are appropriate for
the intended use: did we build the right thing? It requires application,
statistical, scientific, and user evidence.
\end{definitionbox}

\begin{definitionbox}{Reproducibility}
\emph{Reproducibility} is the ability to reconstruct a result under documented
conditions using identified data, code, configuration, dependencies, and
procedures. It controls computational history; it does not prove validity.
\end{definitionbox}

A reproducible error remains an error. A validated claim that cannot be
reconstructed remains difficult to audit or extend. Strong work seeks all
three.

\subsection{Composition and numerical judgment}

Suppose a producer guarantees a temperature in Celsius and a consumer assumes
Celsius. Their composition is meaningful only if the guarantee is true and at
least as strong as the assumption. The type \code{float} cannot express unit,
calibration, valid range, sample identity, or privacy status. Domain classes,
schemas, validation, documentation, and tests make more of that contract
explicit.

Numerical claims also require tolerances. Floating-point equality may be the
wrong contract; an acceptable tolerance must follow from the algorithm, scale,
accumulated error, and downstream use. For iterative methods, report convergence
evidence and nonconvergence. For randomized procedures, a fixed seed aids
diagnosis but evaluation across seeds measures sensitivity.

\begin{workedexample}{High accuracy, invalid scientific evidence}
A laboratory records ten measurements from each of one hundred specimens and
randomly splits the one thousand rows into training and test sets. A classifier
reports \(94\%\) accuracy.

The implementation can be correct and reproducible while the evaluation is
invalid for new specimens: measurements from the same specimen occur in both
sets. If specimen-specific signatures are easy to learn, the score overstates
generalization. Splitting by specimen identity repairs the statistical
contract. A regression test can then require disjoint specimen identifiers,
but domain reasoning is what identified the correct grouping.
\end{workedexample}

\begin{professionalbox}
Report only the precision supported by measurement and computation.
\code{97.2000000000} does not make an instrument more precise. Measurement
uncertainty, numerical error, rounding, and display formatting are distinct.
\end{professionalbox}

\subsection{Failure investigation}

When a result is disputed:

\begin{enumerate}
  \item state the claim, version, conditions, units, tolerance, and intended use;
  \item preserve and reconstruct the inputs, artifact, environment, and output;
  \item trace provenance backward to the first violated contract;
  \item reduce the violation to the smallest representative case;
  \item repair the data, contract, implementation, or study design at its source;
  \item add regression evidence at that layer and affected integrations; and
  \item re-evaluate downstream claims before releasing the correction.
\end{enumerate}

The required evidence is proportional to consequence, but every consequential
claim needs an owner and an evidence portfolio spanning the relevant lenses.

\conceptcheck{A pipeline is reproducible, all tests pass, and the deployed model
matches the evaluated version. Give one statistical and one scientific reason
its conclusion could still be wrong.}

\section{A professional-practice spine}

\begin{definitionbox}{Professional practice}
\emph{Professional practice} is the discipline that makes technical work
understandable, reviewable, reproducible, safe to change, and recoverable. It
coordinates software artifacts with explicit human responsibility.
\end{definitionbox}

The operating loop is:

\begin{center}
\begin{tikzpicture}[node distance=6mm and 5mm]
  \node[flowbox,text width=25mm] (clarify) {Clarify\\claim, owner};
  \node[flowbox,text width=25mm,right=of clarify] (preserve) {Preserve\\evidence, context};
  \node[flowbox,text width=25mm,right=of preserve] (encode) {Encode\\contract, behavior};
  \node[flowbox,text width=25mm,right=of encode] (verify) {Verify\\tests, review};
  \node[flowbox,text width=25mm,below=of verify] (release) {Release\\version, approval};
  \node[flowbox,text width=25mm,left=of release] (observe) {Observe\\signals, outcomes};
  \node[flowbox,text width=25mm,left=of observe] (learn) {Learn\\repair, revision};
  \draw[flowarrow] (clarify) -- (preserve);
  \draw[flowarrow] (preserve) -- (encode);
  \draw[flowarrow] (encode) -- (verify);
  \draw[flowarrow] (verify) -- (release);
  \draw[flowarrow] (release) -- (observe);
  \draw[flowarrow] (observe) -- (learn);
  \draw[flowarrow] (learn) -- (clarify);
\end{tikzpicture}
\end{center}

Each stage has a concrete obligation.

\begin{longtable}{@{}
  >{\raggedright\arraybackslash}p{0.22\textwidth}
  >{\raggedright\arraybackslash}p{0.38\textwidth}
  >{\raggedright\arraybackslash}p{0.31\textwidth}@{}}
\toprule
\textbf{Obligation} & \textbf{Practice} & \textbf{Representative artifact} \\
\midrule
\endhead
Name the boundary & Specify producer, consumer, input, output, failure, and
owner. & Type, schema, interface contract \\
Preserve provenance & Retain source evidence and record transformations without
assuming raw means correct. & Immutable input, checksum, lineage record \\
Make contracts executable & Encode what software can check; document what
requires domain judgment. & Validation, types, tests, database constraint \\
Fail recoverably & Reject invalid states near entry, add safe context, and
retain a known-good state. & Exception contract, rollback procedure \\
Version behavior & Identify code, data, dependencies, model, configuration,
schema, and prompt when they affect results. & Commit, lockfile, dataset and
model versions \\
Bound resources and authority & Limit rows, memory, time, retries, cost,
permissions, and side effects. & Query limit, timeout, scoped credential,
approval gate \\
Observe outcomes & Collect signals that answer defined questions while
protecting secrets and privacy. & Logs, metrics, traces, audits, outcome checks \\
Design for review & Keep changes coherent and state intent, evidence, risk,
compatibility, and recovery. & Pull request, CI result, release record \\
\bottomrule
\end{longtable}

An exception should identify the violated contract and support a safe response.
Catch it only to add context, recover, translate it to an interface response, or
release resources. A retry is appropriate only for a plausibly transient
failure and an operation safe to repeat.

Version control must extend beyond source code. A commit identifies code but
not an external dataset, a moving provider default, or an unrecorded training
configuration. Credentials are configuration but never repository content;
store them in environment variables or a secret service with the narrowest
useful permissions.

Observability is similarly scoped. Logs describe events, metrics summarize
quantities over time, traces follow work across components, and audits record
consequential authority. None is useful without a question, retention policy,
and owner responsible for action.

\begin{misconceptionbox}
\textbf{Misconception:} professional engineering begins only after exploration
succeeds.

\textbf{Correction:} rigor should match consequence, but a named environment,
preserved input, explicit function, focused test, and recorded result make even
exploratory work easier to evaluate and revisit.
\end{misconceptionbox}

\subsection{Scale rigor with consequence}

\begin{center}
\small
\begin{tabularx}{\textwidth}{@{}
  >{\raggedright\arraybackslash}p{0.24\textwidth}
  >{\raggedright\arraybackslash}X
  >{\raggedright\arraybackslash}X@{}}
\toprule
\textbf{Context} & \textbf{Reasonable minimum} & \textbf{Additional controls} \\
\midrule
Private exploration & Named environment, preserved input, recorded assumptions,
reproducible notebook or script & Tests for reused logic, versioned output,
resource bounds \\
Shared analysis & Locked environment, provenance, review, statistical evidence &
Independent reproduction, access control, release record \\
User-facing product & Contract and end-to-end tests, versioned deployment,
monitoring, owner, rollback & Security review, incident response, staged
rollout, client evidence \\
Consequential automation & Explicit authority, evaluation, approval, audit
trail, failure analysis & Formal governance, adversarial testing, continuous
outcome evaluation \\
\bottomrule
\end{tabularx}
\end{center}

\emph{Done} is therefore contextual. At minimum, the intended behavior is
implemented; relevant contracts and versions are identified; the evidence is
appropriate to the claim; and plausible failures have detection, ownership,
and recovery.

\conceptcheck{Why can adding automation reduce engineering quality when
authority, evidence, ownership, and recovery remain implicit?}

\section{How to read and use this book}

The text states concepts and contracts; notebooks make them executable; package
source holds reusable behavior; tests preserve selected expectations. Read by
moving deliberately among these artifacts.

\begin{definitionbox}{Active computational reading}
\emph{Active computational reading} treats prose and code as claims to be
examined. The reader predicts behavior, executes in a known environment,
inspects the result, explains discrepancies, and records reusable evidence.
\end{definitionbox}

\subsection{The predict-execute-explain loop}

For each substantial example:

\begin{enumerate}
  \item state the question and the contract under examination;
  \item predict the value, type, shape, unit, mutation, or exception;
  \item execute the smallest relevant case in a known environment;
  \item inspect values and metadata rather than merely noting success;
  \item explain the observation from the language rule, algorithm, data, and
  environment; and
  \item vary one condition and convert a durable conclusion into a test or note.
\end{enumerate}

Prediction exposes the current mental model. Variation distinguishes an
explanation from a story invented after seeing one output.

\subsection{Read code as a contract}

Ask three questions before execution:

\begin{description}
  \item[Interface.] What enters, returns, or is written? Which types, shapes,
  units, schemas, and exceptions define the boundary?
  \item[Mechanism.] Which operations transform the input? What mutable or
  external state, randomness, file, service, or database affects the result?
  \item[Evidence.] Which ordinary, boundary, and failure cases are exercised?
  What important claim remains untested?
\end{description}

\begin{workedexample}{Interrogating a probability function}
\begin{lstlisting}
def mean_probability(values: list[float]) -> float:
    if not values:
        raise ValueError("values must not be empty")
    if any(value < 0 or value > 1 for value in values):
        raise ValueError("probabilities must lie in [0, 1]")
    return sum(values) / len(values)
\end{lstlisting}

The contract is a nonempty sequence of probabilities mapped to its arithmetic
mean; empty or out-of-range input fails. Predict \code{[0.2, 0.4, 0.6]},
\code{[]}, and \code{[0.4, 1.2]}. Then inspect what the hints do not enforce at
runtime and test \code{nan}. Even complete software tests would not establish
that averaging probabilities is scientifically appropriate.
\end{workedexample}

\subsection{Use notebooks as laboratories}

\begin{definitionbox}{Notebook state}
\emph{Notebook state} is the set of names, imported modules, objects, random
states, environment settings, files, and external resources available to the
kernel. Visible cell order does not prove execution order.
\end{definitionbox}

Before treating a notebook as evidence:

\begin{enumerate}
  \item select and verify the project kernel;
  \item restart it to remove hidden state;
  \item run all cells from top to bottom;
  \item inspect warnings, paths, seeds, inputs, execution order, and invariants;
  \item after the final edit, restart and run all again.
\end{enumerate}

\begin{professionalbox}
A notebook is an experimental laboratory, not an operational boundary. It is
well suited to posing a question, inspecting data, varying assumptions, and
joining evidence to interpretation. A notebook may demonstrate an end-to-end
path without being the end-to-end system. If work must run unattended, serve
concurrent users, retain durable state, enforce access policy, emit operational
telemetry, or support rollback, move that responsibility into a tested package,
script, service, database, or workflow.

For sustained academic research, the notebook may remain the interpretive
record, but it must not be the only source of truth. Data identity, provenance,
configuration, transformations, and reusable computation must be reconstructible
outside hidden kernel state.
\end{professionalbox}

When behavior differs from the text, preserve the error, read the final
exception line, locate the first frame in code you control, and reduce the
problem before changing dependencies. For \code{ModuleNotFoundError}, first
compare \code{sys.executable} with the project environment. For a missing file,
inspect the working directory and use \code{pathlib.Path}. For a numerical
difference, record inputs, versions, dtype, shape, seed, and tolerance.

\subsection{Move reusable work into durable artifacts}

An exploratory expression belongs in a notebook. Repeated or consequential
logic belongs in a named function with type hints, a NumPy-style docstring, and
explicit failure behavior. Package source provides a stable import boundary;
tests preserve ordinary, edge, and invalid cases; version history records the
change and its rationale. Import the tested function back into the notebook so
that the notebook carries the analytical argument rather than a private
implementation.

Use a broader graduation rule for responsibilities that are not merely Python
functions: repeatable batch work moves to a script or command-line interface;
shared state moves to durable storage; concurrent access moves behind a service;
and scheduling, deployment, monitoring, and rollback move to operational
workflows. Extract a cell when its behavior must survive reuse, automation,
collaboration, publication, or failure.

A short computational record should identify the question, code and data
versions, environment, prediction, observed evidence, interpretation, and next
action. The standard is reconstruction, not length.

\conceptcheck{Why is a cleanly executed notebook necessary but insufficient
evidence that its author understands the computation or that its scientific
claim is valid?}

\section{A running case study}

Part I repeatedly enlarges a scientific knowledge system. It begins with small
records describing entities and relationships:

\begin{enumerate}
  \item native Python values establish identifiers, labels, and edges;
  \item classes state invariants, and functions reject invalid relationships;
  \item CSV and JSON preserve records, with round-trip tests for identity;
  \item a package and CI make transformations reusable and reviewable;
  \item arrays, tables, and figures support numerical investigation;
  \item a database retrieves bounded neighborhoods without loading everything;
  \item a language-model tool may propose descriptions or edges under schema
  validation and human review; and
  \item an API exposes validated queries to a client with operational evidence.
\end{enumerate}

The point is cumulative design: the deployed product still depends on the
meaning of a dictionary key and the exception contract defined much earlier.
Examples will vary across probability, mathematics, physics, chemistry, and
industrial engineering so that the software principles do not depend on one
domain.

\conceptcheck{For one value in this system, list the identity, unit or schema,
provenance, transformation, version, and uncertainty required to interpret it.}

\newpage
\section{Exercises}

\subsection*{Foundational}
\begin{enumerate}
  \item Distinguish a model, a data product, and a data-driven software system.
  Give one example in which all three have different owners.
  \item For the claim \emph{the reported temperature is \(97.2\)}, list the
  missing context under each of the six correctness lenses.
  \item Explain why verification, validation, and reproducibility are mutually
  supportive but logically distinct.
\end{enumerate}

\subsection*{Applied}
Select one result from a class, laboratory, or public dataset. Trace it backward
to source evidence and forward to its audience. At each boundary, record the
representation, contract, owner, likely failure, and available evidence.

\subsection*{Design}
Write a one-page system charter for a small data product. Specify its question,
users, evidence, output contract, privacy and safety constraints, version
identity, resource and authority bounds, monitoring, and recovery. Choose
technologies only after these responsibilities are explicit.

\begin{chapterreview}
\textbf{Key ideas.} A learned model is one component of an evidence-bearing
software system. Correctness claims are scoped by assumptions and require
different evidence at data, computational, software, statistical, scientific,
and operational layers. Professional practice makes boundaries, versions,
authority, ownership, observation, and recovery explicit.

\textbf{Vocabulary.} model; data product; boundary; contract; provenance;
verification; validation; reproducibility; observability; responsibility.

\textbf{Explain aloud.} How can a model with excellent benchmark performance
belong to an invalid or unsafe data product?
\end{chapterreview}

\begin{notebooklink}
The operational checklist is in
\code{notebooks/PROFESSIONAL\_PRACTICES.md}. Lecture notebooks supply the
executable investigations.
\end{notebooklink}
